%============================================================
%  Bd. 17 - Idempotente Halbgruppen
%============================================================

\documentclass[main.tex]{subfiles}

\ifSubfilesClassLoaded{
  \BandSetupStandalone{B17}
}{
}

\title{Bd. 17 - Idempotente Halbgruppen}
\author{Martin Kunze}
\date{}
\setcounter{file}{17}

\begin{document}

\title{Bd. 17 - Idempotente Halbgruppen}
\author{Martin Kunze}
\date{}
\hypersetup{pageanchor=false}
\maketitle
\hypersetup{pageanchor=true}
\setcounter{tocdepth}{3}
\tableofcontents
\setcounter{file}{17}
\renewcommand{\FormulaBandID}{17}

\chapter{Einführung}

Dieser Band untersucht Halbgruppen mit idempotenter Operation.

\chapter{Idempotente Halbgruppen}

\section{Definition}

\FormulaDefDeltaK[Begriff der idempotenten Halbgruppe]{%
  \IdempotentSemigroupAlg{A}{\star}%
}{IdempotentSemigroupDef}{%
  \DeltaRow{Mengen}{A\dsep x}
  \DeltaRow{\textbf{Axiome}}{}
  \DeltaPrem{Halbgruppe}{\SemigroupAlg{A}{\star}}
  \DeltaRow{Idempotenz}{x\in A\vdash x\star x=x}
  \DeltaRow{\textbf{Neues Symbol}}{}
  \DeltaPrem{Idempotente Halbgruppe}{%
    \IdempotentSemigroupAlg{A}{\star}}
}

\FormulaAxiomDeltaK[Zugriff auf die Halbgruppe einer idempotenten Halbgruppe]{%
  \IdempotentSemigroupAlg{A}{\star}
  \vdash\SemigroupAlg{A}{\star}%
}{IdempotentSemigroupBaseAxiom}{%
  \DeltaRow{Mengen}{A}
}

\FormulaAxiomDeltaK[Zugriff auf die Idempotenz]{%
  \IdempotentSemigroupAlg{A}{\star}\dsep x\in A
  \vdash x\star x=x%
}{IdempotentSemigroupIdempotenceAxiom}{%
  \DeltaRow{Mengen}{A\dsep x}
}

\section{Beispiele}

\FormulaThmDeltaK[Die Vereinigung bildet auf einer Potenzmenge eine
idempotente Halbgruppe]{%
  \IdempotentSemigroupAlg{\powerset(M)}{\cup}%
}{PowerSetUnionIdempotentSemigroup}{%
  \DeltaRow{Mengen}{M\dsep X}
}
\begin{tabproof}
  \proofstep{}{\CommutativeSemigroupAlg{\powerset(M)}{\cup}}{%
    \FormulaRefAuto{PowerSetUnionCommutativeSemigroup}[thm]}
  \proofstep{}{\SemigroupAlg{\powerset(M)}{\cup}}{%
    \FormulaRefAuto{CommutativeSemigroupBaseAxiom}[ax]{1}}
  \proofstep{3}{X\in\powerset(M)}{\rA}
  \proofstep{}{X\cup X=X}{%
    \FormulaRefAuto{A \cup A = A}}
  \proofstep{}{\IdempotentSemigroupAlg{\powerset(M)}{\cup}}{%
    \FormulaRefAuto{IdempotentSemigroupDef}[def]{2,4}}
\end{tabproof}

\paragraph{Nichtkommutatives Beispiel.}
Für eine Menge $A$ mit mindestens zwei Elementen sei
$x\star_{\ell}y\coloneqq x$. Dann gelten
$(x\star_{\ell}y)\star_{\ell}z=x
=x\star_{\ell}(y\star_{\ell}z)$ und
$x\star_{\ell}x=x$, aber im Allgemeinen
$x\star_{\ell}y\neq y\star_{\ell}x$.

\section{Morphismen}

\FormulaDefDeltaK[Homomorphismus idempotenter Halbgruppen]{%
  \IdempotentSemigroupHom{f}{A,\star}{B,\diamond}%
}{IdempotentSemigroupHomDef}{%
  \DeltaRow{Mengen}{A\dsep B}
  \DeltaRow{Funktionen}{f\dsep\star\dsep\diamond}
  \DeltaRow{\textbf{Axiome}}{}
  \DeltaPrem{Quellstruktur}{\IdempotentSemigroupAlg{A}{\star}}
  \DeltaPrem{Zielstruktur}{\IdempotentSemigroupAlg{B}{\diamond}}
  \DeltaPrem{Halbgruppenhomomorphismus}{%
    \SemigroupHom{f}{A,\star}{B,\diamond}}
  \DeltaRow{\textbf{Neues Symbol}}{}
  \DeltaPrem{Homomorphismus}{%
    \IdempotentSemigroupHom{f}{A,\star}{B,\diamond}}
}

\FormulaAxiomDeltaK[Quellstruktur]{%
  \IdempotentSemigroupHom{f}{A,\star}{B,\diamond}
  \vdash\IdempotentSemigroupAlg{A}{\star}%
}{IdempotentSemigroupHomSourceAxiom}{%
  \DeltaRow{Mengen}{A\dsep B}
  \DeltaRow{Funktionen}{f\dsep\star\dsep\diamond}
}

\FormulaAxiomDeltaK[Zielstruktur]{%
  \IdempotentSemigroupHom{f}{A,\star}{B,\diamond}
  \vdash\IdempotentSemigroupAlg{B}{\diamond}%
}{IdempotentSemigroupHomTargetAxiom}{%
  \DeltaRow{Mengen}{A\dsep B}
  \DeltaRow{Funktionen}{f\dsep\star\dsep\diamond}
}

\FormulaAxiomDeltaK[Zugriff auf den Halbgruppenhomomorphismus]{%
  \IdempotentSemigroupHom{f}{A,\star}{B,\diamond}
  \vdash\SemigroupHom{f}{A,\star}{B,\diamond}%
}{IdempotentSemigroupHomBaseHomAxiom}{%
  \DeltaRow{Mengen}{A\dsep B}
  \DeltaRow{Funktionen}{f\dsep\star\dsep\diamond}
}

\FormulaDefDeltaK[Isomorphismus idempotenter Halbgruppen]{%
  \IdempotentSemigroupIso{f}{A,\star}{B,\diamond}%
}{IdempotentSemigroupIsoDef}{%
  \DeltaRow{Mengen}{A\dsep B}
  \DeltaRow{Funktionen}{f\dsep\star\dsep\diamond}
  \DeltaRow{\textbf{Axiome}}{}
  \DeltaPrem{Homomorphismus}{%
    \IdempotentSemigroupHom{f}{A,\star}{B,\diamond}}
  \DeltaPrem{Bijektivität}{f\colon A\bij B}
  \DeltaRow{\textbf{Neues Symbol}}{}
  \DeltaPrem{Isomorphismus}{%
    \IdempotentSemigroupIso{f}{A,\star}{B,\diamond}}
}

\FormulaAxiomDeltaK[Homomorphismuszugriff]{%
  \IdempotentSemigroupIso{f}{A,\star}{B,\diamond}
  \vdash\IdempotentSemigroupHom{f}{A,\star}{B,\diamond}%
}{IdempotentSemigroupIsoHomAxiom}{%
  \DeltaRow{Mengen}{A\dsep B}
  \DeltaRow{Funktionen}{f\dsep\star\dsep\diamond}
}

\FormulaAxiomDeltaK[Bijektivitätszugriff]{%
  \IdempotentSemigroupIso{f}{A,\star}{B,\diamond}
  \vdash f\colon A\bij B%
}{IdempotentSemigroupIsoBijectiveAxiom}{%
  \DeltaRow{Mengen}{A\dsep B}
  \DeltaRow{Funktionen}{f\dsep\star\dsep\diamond}
}

\end{document}
